\begin{exercisechunk}
\item \textbf{From noise-driven maps to an HMM law.}
\label[exercise]{ex:l5-kernels}
Use \cref{eq:nonlinear-ssm,eq:ssm-joint} and the finite HMM matrices
\(A,C\) defined in \cref{eq:hmm-joint}.
\begin{enumerate}[(a),beginpenalty=10000]
\item Prove the joint-law product formula by successive conditioning.
Write the marginal law of \(S_{1:T},Z_{1:T}\) after integrating out \(S_0\).
\item For independent uniform \([0,1]\) noises, give explicit
inverse-cumulative-sum maps implementing \(A\) and \(C\), including a
convention at zero-probability endpoints.
\item Construct two different noise maps with the same transition kernel.
Explain why a specified model gives a unique law but the law does not
uniquely specify its noise maps.
\item Suppose the hidden states are observed. Derive the column-normalized
transition and emission count estimates. Which columns remain undetermined
if some states are never visited?
\end{enumerate}
\end{exercisechunk}
